% Bilingual EN/KO proposal deck on the UBI beamer theme.
%
% Build (from this directory, twice so the frame count resolves):
%     lualatex presentacion.tex
%
% LuaLaTeX, not XeLaTeX. Under XeLaTeX this theme renders only its first two
% frames and every later one comes out blank; under LuaLaTeX — the engine its
% README was written for — it is fine. Korean therefore comes from a luaotfload
% font fallback rather than from xeCJK.
%
% 4:3 is the theme's native ratio; its background art and text block are drawn
% for it.
%
% The descriptive table comes from ../tables_naver.tex, generated by
% 10_01_descriptive_tables.py — re-run that script and recompile to refresh it.
%
% Scope: Naver Place only.

\documentclass[compress]{beamer}

\usetheme[dark,framenumber,totalframenumber]{UBI}

\usepackage{fontspec}
\usepackage{microtype}
% Georgia is the UBI house font and carries no Hangul, so Korean falls back to
% Malgun Gothic glyph by glyph. This keeps Latin text in Georgia and leaves the
% spaces between Korean words intact.
\directlua{luaotfload.add_fallback("hangul", {"Malgun Gothic:mode=node;"})}
\setmainfont{Georgia}[RawFeature={fallback=hangul}]
\setsansfont{Georgia}[RawFeature={fallback=hangul}]

\usepackage{booktabs}
\usepackage{amsmath}
\usepackage{float}

% The theme sets `normal text` to UBIgrey (RGB 191,191,191) — legible on the
% dark title slide, all but invisible on the light content frames.
\setbeamercolor{normal text}{fg=black!87}
\setbeamercolor{itemize item}{fg=UBIblue}
\setbeamercolor{enumerate item}{fg=UBIblue}

\setbeamertemplate{caption}{\insertcaption}
\setbeamerfont{frametitle}{size=\large}
\setbeamerfont{title}{size=\Large}
\setbeamerfont{subtitle}{size=\normalsize}

% Korean gloss: beside a frame title, and under running text.
\newcommand{\ko}[1]{{\footnotesize\color{black!50}#1}}
\newcommand{\kosub}[1]{{\scriptsize\color{black!55}#1}}
% Boxed aside, in the theme's own blue.
\newcommand{\aside}[1]{%
  \begin{center}\colorbox{UBIblue!12}{%
    \parbox{\dimexpr\textwidth-14pt\relax}{\vspace{2pt}#1\vspace{2pt}}}\end{center}}

% Generated by 10_01_descriptive_tables.py -- do not edit by hand.
% Requires \usepackage{float} in the including document ([H] placement).
% Defines: \tblPlaces \tblVariables \tblDescriptives \tblCorrelations \tblVIF

\newcommand{\tblPlaces}{%
\begin{table}[H]\centering%
\caption{관광지별 리뷰 수 / Reviews per attraction}%
\footnotesize\begin{tabular}{llrrr}\hline%
관광지 & English & Reviews & w/ text & w/ rating \\ \hline%
서울스카이 & Seoul Sky & 1,092 & 1,025 & 792 \\%
남산골한옥마을 & Namsangol Village & 1,075 & 1,033 & 826 \\%
에버랜드 & Everland & 995 & 935 & 851 \\%
N서울타워 & N Seoul Tower & 939 & 817 & 198 \\%
동대문디자인플라자 & DDP & 668 & 453 & 232 \\%
덕수궁 & Deoksugung & 646 & 539 & 72 \\%
운현궁 & Unhyeongung & 616 & 562 & 384 \\%
서대문형무소역사관 & Seodaemun Prison History Hall & 610 & 331 & 27 \\%
창덕궁 & Changdeokgung & 610 & 399 & 38 \\%
올림픽공원 & Olympic Park & 609 & 499 & 104 \\%
해운대해수욕장 & Haeundae Beach & 608 & 457 & 53 \\%
서대문자연사박물관 & Seodaemun Museum & 608 & 320 & 20 \\%
광안리해수욕장 & Gwangalli Beach & 607 & 505 & 77 \\%
국립중앙박물관 & Natl. Museum & 606 & 543 & 221 \\%
여의도한강공원 & Yeouido Park & 605 & 443 & 87 \\%
코엑스 & COEX & 590 & 359 & 198 \\%
창경궁 & Changgyeonggung & 590 & 418 & 26 \\%
경복궁 & Gyeongbokgung & 580 & 477 & 91 \\%
낙산공원 & Naksan Park & 435 & 249 & 108 \\%
북촌한옥마을 & Bukchon Village & 296 & 165 & 118 \\%
북한산 & Bukhansan & 253 & 155 & 26 \\%
롯데월드 & Lotte World & 180 & 85 & 10 \\%
경희궁 & Gyeonghuigung & 170 & 98 & 12 \\%
은평한옥마을 & Eunpyeong Village & 156 & 125 & 36 \\%
인사동길 & Insadong-gil & 109 & 68 & 6 \\%
광장시장 & Gwangjang Mkt & 59 & 32 & 14 \\%
명동성당 & Myeongdong Cathedral & 39 & 33 & 18 \\%
홍대거리 & Hongdae St. & 30 & 22 & 6 \\%
서촌 & Seochon & 20 & 18 & 1 \\%
\hline\end{tabular}%
\end{table}%
}

\newcommand{\tblVariables}{%
\begin{table}[H]\centering%
\caption{변수요약 / Variable summary}%
\footnotesize\begin{tabular}{llp{7.6cm}}\hline%
변수 종류 & 변수 이름 & 변수 설명 \\ \hline%
종속변수 (DV) & Revisit (재방문) & 해당 관광지를 2회차 이상 방문한 작성자의 리뷰 여부 (0/1) \\%
 & Keyword Count (키워드 수) & 리뷰당 작성자가 선택한 키워드 수 (0$\sim$5), 정보 다양성 지표 \\%
 & Rating (별점) & 작성자가 부여한 별점 (1$\sim$5) \\%
 & Reactions (반응 수) & 해당 리뷰에 반응을 남긴 이용자 수, 유용성 지표 \\%
독립변수 (IV) & Companion (동반자) & 방문 동반자 유형 더미: 혼자·연인·친구·아이·가족·동료 \\%
 & Purpose (방문 목적) & 방문 목적 더미: 여행·나들이·일상·데이트·친목 \\%
 & Waited (대기) & 방문 시 대기 경험 여부 (0/1) \\%
 & Reserved (예약) & 사전 예약 여부 (0/1) \\%
 & Sentiment (감성 점수) & 리뷰 텍스트의 감성 점수 (-1 부정 $\sim$ 1 긍정) \\%
 & Tag themes (키워드 주제) & 키워드를 10개 주제로 분류한 더미: 시각·볼거리·혼잡·한적·시설·접근성·가족·가격·음식·물놀이 \\%
통제변수 (Control) & Text Length (텍스트 길이) & 리뷰 텍스트 길이의 로그값 \\%
 & Image Count (사진 수) & 리뷰에 첨부된 사진 수 \\%
 & User Reviews (작성자 리뷰 수) & 작성자가 누적 작성한 리뷰 수의 로그값 \\%
 & User Followers (팔로워 수) & 작성자의 팔로워 수의 로그값 \\%
 & Review Age (리뷰 경과일) & 리뷰 작성 후 경과일의 로그값, 반응 수 모형에서 노출 기간을 통제 \\%
 & Verification (인증) & 영수증 인증 리뷰 여부 (0/1) \\%
 & Season (계절) & 방문 계절 더미: 여름·가을·겨울 (봄 기준) \\%
 & Weekend (주말) & 주말 방문 여부 (0/1) \\%
 & Place FE (관광지 고정효과) & 관광지별 고정효과, 장소 고유 특성 통제 \\%
\hline\end{tabular}\end{table}%
}

\newcommand{\tblDescriptives}{%
\begin{table}[H]\centering%
\caption{기술통계표 / Descriptive statistics}%
\footnotesize\begin{tabular}{lrrrrrr}\hline%
Statistic & N & Mean & St.\ Dev. & Min & Median & Max \\ \hline%
Revisit (재방문) & 14,401 & 0.218 & 0.413 & 0 & 0 & 1 \\%
Keyword Count (키워드 수) & 14,401 & 2.97 & 1.883 & 0 & 3 & 5 \\%
Rating (별점) & 4,652 & 4.682 & 0.732 & 0.5 & 5 & 5 \\%
Reactions (반응 수) & 14,401 & 1.437 & 7.415 & 0 & 0 & 122 \\%
Sentiment (감성 점수) & 11,165 & 0.653 & 0.457 & -0.973 & 0.87 & 0.988 \\%
Text Length (텍스트 길이) & 11,165 & 112.089 & 106.805 & 14 & 71 & 753 \\%
Image Count (사진 수) & 14,401 & 0.235 & 1.412 & 0 & 0 & 20 \\%
User Reviews (작성자 리뷰 수) & 14,401 & 823.26 & 1,009.4 & 1 & 482 & 19,082.0 \\%
Review Age (리뷰 경과일) & 14,401 & 588.324 & 558.142 & 0 & 428 & 3,957.0 \\%
\hline\end{tabular}\end{table}%
}

\newcommand{\tblCorrelations}{%
\begin{table}[H]\centering%
\caption{상관분석표 / Correlation matrix}%
\footnotesize\begin{tabular}{lrrrrrrrrr}\hline%
 & 1 & 2 & 3 & 4 & 5 & 6 & 7 & 8 & 9 \\ \hline%
1. Revisit (재방문) & 1.000 &  &  &  &  &  &  &  &  \\%
2. Keyword Count (키워드 수) & 0.140 & 1.000 &  &  &  &  &  &  &  \\%
3. Rating (별점) & 0.069 & 0.264 & 1.000 &  &  &  &  &  &  \\%
4. Reactions (반응 수) & 0.194 & 0.097 & 0.051 & 1.000 &  &  &  &  &  \\%
5. Sentiment (감성 점수) & -0.067 & 0.116 & 0.514 & -0.039 & 1.000 &  &  &  &  \\%
6. Text Length (텍스트 길이) & 0.198 & 0.204 & 0.000 & 0.221 & -0.095 & 1.000 &  &  &  \\%
7. Image Count (사진 수) & 0.036 & 0.075 & 0.044 & 0.057 & 0.016 & 0.101 & 1.000 &  &  \\%
8. User Reviews (작성자 리뷰 수) & 0.255 & 0.031 & 0.036 & 0.212 & -0.082 & 0.159 & 0.046 & 1.000 &  \\%
9. Review Age (리뷰 경과일) & -0.140 & -0.395 & -0.094 & -0.083 & -0.027 & -0.137 & -0.057 & 0.034 & 1.000 \\%
\hline\end{tabular}\end{table}%
}

\newcommand{\tblVIF}{%
\begin{table}[H]\centering%
\caption{분산팽창계수 / Variance inflation factors}%
\footnotesize\begin{tabular}{lrlrlr}\hline%
변수 / Variable & VIF & 변수 / Variable & VIF & 변수 / Variable & VIF \\ \hline%
has\_visit\_info & 2.82 & comp\_family & 1.21 & sentiment & 1.06 \\%
waited & 1.10 & comp\_colleague & 1.25 & n\_tags & 1.54 \\%
reserved & 1.84 & purp\_travel & 1.38 & revisit & 1.16 \\%
comp\_alone & 1.69 & purp\_outing & 1.65 & ln\_length & 1.10 \\%
comp\_partner & 1.75 & purp\_daily & 1.45 & ln\_user\_reviews & 1.54 \\%
comp\_friend & 1.53 & purp\_date & 1.67 &  &  \\%
comp\_child & 1.70 & purp\_social & 1.39 &  &  \\%
\hline\end{tabular}%
\par\vspace{2pt}\footnotesize\textit{All values below 5; multicollinearity is not a concern. 모든 값이 5 미만으로 다중공선성 문제는 없다.}%
\end{table}%
}

% Shared body for the variable-map table; used by \tblVariableMapEN and \tblVariableMapKO.
% Notation key: $Y_{ij}$ = outcome of that model; $\beta_k$ = coefficient in x_ij vector;
%   $\delta_j$ = place fixed effect; $\gamma$ = ln_review_age exposure offset (M4 only);
%   $\beta_k{\times}N$ = N binary dummies each with their own \beta_k; -- = not included.
% Scale abbreviations used below:
%   bin. = binary (0/1)   ord. = ordinal   cont. = continuous   count = count (integer)   categ. = categorical
\newcommand{\tblVariableMapEN}{%
\begin{table}[H]\centering%
\caption{Variable-to-model mapping}%
\footnotesize\setlength{\tabcolsep}{3.5pt}%
\begin{tabular}{lp{2.7cm}p{2.8cm}lp{1.4cm}p{1.4cm}p{1.4cm}p{1.4cm}}\hline%
Type & Concept (Table~2) & Dataset column(s) & Scale & M1 & M2 & M3 & M4 \\\hline%
\multicolumn{8}{l}{\quad\textit{Dependent variables}}\\\hline%
DV & Revisit             & \texttt{revisit}                          & bin.          & $Y_{ij}$              & $\beta_k$             & $\beta_k$             & $\beta_k$             \\%
DV & Keyword Count       & \texttt{n\_tags}                          & ord.\ 0--5    & $\beta_k$             & $Y_{ij}$              & $\beta_k$             & $\beta_k$             \\%
DV & Rating              & \texttt{rating}                           & ord.\ 1--5    & --                    & --                    & $Y_{ij}$              & --                    \\%
DV & Reactions           & \texttt{reactions}                        & count         & --                    & --                    & --                    & $Y_{ij}$              \\\hline%
\multicolumn{8}{l}{\quad\textit{Independent variables}}\\\hline%
IV & Companion           & \texttt{comp\_*} ($\times$6)              & bin.          & $\beta_k{\times}6$    & $\beta_k{\times}6$    & $\beta_k{\times}4^a$  & $\beta_k{\times}6$    \\%
IV & Purpose             & \texttt{purp\_*} ($\times$5)              & bin.          & $\beta_k{\times}5$    & $\beta_k{\times}5$    & $\beta_k{\times}2^a$  & $\beta_k{\times}5$    \\%
IV & Waited              & \texttt{waited}                           & bin.          & $\beta_k$             & $\beta_k$             & --                    & $\beta_k$             \\%
IV & Reserved            & \texttt{reserved}                         & bin.          & $\beta_k$             & $\beta_k$             & $\beta_k$             & $\beta_k$             \\%
IV & Sentiment           & \texttt{sentiment}                        & cont.         & $\beta_k$             & --                    & --$^b$                & $\beta_k$             \\%
IV & Tag themes          & \texttt{tag\_*} ($\times$10)              & bin.          & --                    & --                    & --                    & --$^c$                \\%
IV & Has text            & \texttt{has\_text}                        & bin.          & --                    & $\beta_k$             & $\beta_k$             & --                    \\%
IV & Visit info          & \texttt{has\_visit\_info}                 & bin.          & $\beta_k$             & $\beta_k$             & --                    & $\beta_k$             \\\hline%
\multicolumn{8}{l}{\quad\textit{Control variables}}\\\hline%
Ctrl & Text Length       & \texttt{ln\_length}                       & cont.         & $\beta_k$             & --                    & $\beta_k$             & $\beta_k$             \\%
Ctrl & User Reviews      & \texttt{ln\_user\_reviews}                & cont.         & $\beta_k$             & $\beta_k$             & $\beta_k$             & $\beta_k$             \\%
Ctrl & User Followers    & \texttt{ln\_user\_followers}              & cont.         & --                    & --                    & --                    & $\beta_k$             \\%
Ctrl & Review Age        & \texttt{ln\_review\_age}                  & cont.         & --                    & --                    & --                    & $\gamma$              \\%
Ctrl & Season            & \texttt{season\_*} ($\times$3)            & bin.          & --                    & --                    & --                    & $\beta_k{\times}3$    \\%
Ctrl & Weekend           & \texttt{visit\_weekend}                   & bin.          & --                    & --                    & --                    & $\beta_k$             \\%
Ctrl & Verification      & \texttt{verif\_receipt}                   & bin.          & --                    & --                    & --                    & $\beta_k$             \\%
Ctrl & Place FE          & \textit{(absorbed)}                       & categ.        & $\delta_j$            & $\delta_j$            & $\delta_j$            & $\delta_j$            \\\hline%
\end{tabular}%
\par\vspace{3pt}\footnotesize%
\textit{Scale:} bin.\ = binary (0/1); ord.\ = ordinal; cont.\ = continuous (log-transformed); count = non-negative integer.\quad
\textit{Notation link:}
$\beta_k$ = enters $\mathbf{x}_{ij}^\top\boldsymbol\beta$ with coefficient $\beta_k$;\quad
$Y_{ij}$ = outcome (left-hand side; not in $\mathbf{x}_{ij}$);\quad
$\delta_j$ = place fixed effect;\quad
$\gamma$ = \texttt{ln\_review\_age} as exposure offset in M4 (not a $\beta_k$);\quad
$\beta_k{\times}N$ = $N$ dummies each with their own $\beta_k$.\\%
$^a$~Partial dummies in M3: only 32\% of reviews have a rating.~%
$^b$~Excluded from M3: $r{=}0.51$ with rating; VIF risk.~%
$^c$~Tag themes only in extended models M1b, M1x.%
\end{table}%
}

\newcommand{\tblVariableMapKO}{%
\begin{table}[H]\centering%
\caption{변수 매핑표}%
\footnotesize\setlength{\tabcolsep}{3.5pt}%
\begin{tabular}{lp{2.7cm}p{2.8cm}lp{1.4cm}p{1.4cm}p{1.4cm}p{1.4cm}}\hline%
유형 & 변수 이름 (표~2) & 데이터셋 열 & 척도 & M1 & M2 & M3 & M4 \\\hline%
\multicolumn{8}{l}{\quad\textit{종속변수}}\\\hline%
종속 & 재방문              & \texttt{revisit}                          & 이진          & $Y_{ij}$              & $\beta_k$             & $\beta_k$             & $\beta_k$             \\%
종속 & 키워드 수           & \texttt{n\_tags}                          & 순서 0--5     & $\beta_k$             & $Y_{ij}$              & $\beta_k$             & $\beta_k$             \\%
종속 & 별점               & \texttt{rating}                           & 순서 1--5     & --                    & --                    & $Y_{ij}$              & --                    \\%
종속 & 반응 수            & \texttt{reactions}                        & 계수          & --                    & --                    & --                    & $Y_{ij}$              \\\hline%
\multicolumn{8}{l}{\quad\textit{독립변수}}\\\hline%
독립 & 동반자 유형         & \texttt{comp\_*} ($\times$6)              & 이진          & $\beta_k{\times}6$    & $\beta_k{\times}6$    & $\beta_k{\times}4^a$  & $\beta_k{\times}6$    \\%
독립 & 방문 목적           & \texttt{purp\_*} ($\times$5)              & 이진          & $\beta_k{\times}5$    & $\beta_k{\times}5$    & $\beta_k{\times}2^a$  & $\beta_k{\times}5$    \\%
독립 & 대기 경험           & \texttt{waited}                           & 이진          & $\beta_k$             & $\beta_k$             & --                    & $\beta_k$             \\%
독립 & 사전 예약           & \texttt{reserved}                         & 이진          & $\beta_k$             & $\beta_k$             & $\beta_k$             & $\beta_k$             \\%
독립 & 감성 점수           & \texttt{sentiment}                        & 연속          & $\beta_k$             & --                    & --$^b$                & $\beta_k$             \\%
독립 & 키워드 주제         & \texttt{tag\_*} ($\times$10)              & 이진          & --                    & --                    & --                    & --$^c$                \\%
독립 & 텍스트 입력 여부    & \texttt{has\_text}                        & 이진          & --                    & $\beta_k$             & $\beta_k$             & --                    \\%
독립 & 방문 정보 제공      & \texttt{has\_visit\_info}                 & 이진          & $\beta_k$             & $\beta_k$             & --                    & $\beta_k$             \\\hline%
\multicolumn{8}{l}{\quad\textit{통제변수}}\\\hline%
통제 & 텍스트 길이         & \texttt{ln\_length}                       & 연속          & $\beta_k$             & --                    & $\beta_k$             & $\beta_k$             \\%
통제 & 작성자 리뷰 수      & \texttt{ln\_user\_reviews}                & 연속          & $\beta_k$             & $\beta_k$             & $\beta_k$             & $\beta_k$             \\%
통제 & 팔로워 수           & \texttt{ln\_user\_followers}              & 연속          & --                    & --                    & --                    & $\beta_k$             \\%
통제 & 리뷰 경과일         & \texttt{ln\_review\_age}                  & 연속          & --                    & --                    & --                    & $\gamma$              \\%
통제 & 방문 계절           & \texttt{season\_*} ($\times$3)            & 이진          & --                    & --                    & --                    & $\beta_k{\times}3$    \\%
통제 & 주말 방문           & \texttt{visit\_weekend}                   & 이진          & --                    & --                    & --                    & $\beta_k$             \\%
통제 & 인증 (영수증)       & \texttt{verif\_receipt}                   & 이진          & --                    & --                    & --                    & $\beta_k$             \\%
통제 & 관광지 고정효과     & \textit{(흡수됨)}                         & 범주          & $\delta_j$            & $\delta_j$            & $\delta_j$            & $\delta_j$            \\\hline%
\end{tabular}%
\par\vspace{3pt}\footnotesize%
\textit{척도:} 이진 = 0/1 더미; 순서 = 순서형 정수; 연속 = 로그 변환 연속형; 계수 = 비음 정수.\quad
\textit{기호 연결:}
$\beta_k$ = $\mathbf{x}_{ij}^\top\boldsymbol\beta$ 벡터에 계수 $\beta_k$로 포함;\quad
$Y_{ij}$ = 종속변수(좌변; $\mathbf{x}_{ij}$에 미포함);\quad
$\delta_j$ = 관광지 고정효과;\quad
$\gamma$ = M4에서 노출 오프셋($\beta_k$ 아님);\quad
$\beta_k{\times}N$ = 더미 $N$개가 각각 고유한 계수로 포함.\\%
$^a$~M3 일부 더미: 별점 보유 리뷰 32\%에 불과.~%
$^b$~M3 제외: 별점과 $r{=}0.51$, 다중공선성 우려.~%
$^c$~키워드 주제는 확장 모형 M1b, M1x에서만.%
\end{table}%
}


\title{What Makes an Attraction Review Informative}
\subtitle{Visit Context, Information Diversity and\\ Reader Response on Naver Place\\[4pt]
          \small 네이버 플레이스의 방문 맥락, 정보 다양성, 이용자 반응}
\author{Jorge Adan Soriano Hernandez\\[3pt]
        \small Department of Data Science, Hoseo University\\
        \small 호서대학교 데이터사이언스학과\\[3pt]
        \small October 2026 \quad 2026년 10월}

% ════════════════════════════════════════════════════════════════════
\begin{document}

\maketitle

% ── 목차 ────────────────────────────────────────────────────────────
\begin{frame}
\frametitle{Table of Contents \ko{목차}}
\small
\begin{enumerate}\itemsep4pt
  \item \textbf{Introduction} \ko{서론}
  \item \textbf{Motivation \& Problem} \ko{동기 및 문제 정의}
  \item \textbf{Related Work} \ko{관련 연구}
  \item \textbf{Methodology} \ko{방법론}
  \item \textbf{Results} \ko{결과}
  \item \textbf{Conclusion} \ko{결론}
\end{enumerate}
\end{frame}

% ════════════════════════════════════════════════════════════════════
\begin{frame}
\frametitle{1. Introduction \ko{서론}}
\small
\begin{itemize}\itemsep6pt
  \item Reviews are what travellers consult before choosing where to go.
  \item Lee, Jeong and Kim (2026): information \textbf{diversity} — distinct
        keywords per review — drives performance, separately from
        \textbf{quantity}.
  \item They measured it by business and week. We take it down to the
        \textbf{individual review}.
\end{itemize}
\vspace{6pt}
\kosub{본 연구는 정보 다양성을 개별 리뷰 수준에서 분석한다.}
\end{frame}

% ════════════════════════════════════════════════════════════════════
\begin{frame}
\frametitle{2. Motivation \ko{동기}}
\small
\begin{itemize}\itemsep8pt
  \item Platforms rank reviews without knowing \textbf{who} supplies varied
        information.
  \item Managers see star ratings — but ratings barely move. The signal is in
        what reviewers \textbf{mention}.
\end{itemize}
\vspace{6pt}
\kosub{별점은 거의 변하지 않는다. 유용한 신호는 작성자가 무엇을 언급하는가에 있다.}
\end{frame}

\begin{frame}
\frametitle{2. Problem Statement \ko{문제 정의}}
\aside{\small
  Lee et al.\ (2026) observe performance, but aggregate to the
  \textbf{business--week} level — the content and visit context of any single
  review stay invisible.\par\vspace{5pt}
  Companion, purpose and waiting time are never used, though Naver collects all
  three.\par\vspace{5pt}
  {\color{UBIred}We know diversity matters — not who supplies it.}}
\vspace{4pt}
\kosub{다양성이 중요하다는 점은 알지만, 누가 그것을 제공하는지는 모른다.}
\end{frame}

% ════════════════════════════════════════════════════════════════════
\begin{frame}
\frametitle{3. Related Work \ko{관련 연구}}
\begin{center}
{\scriptsize\color{black}
\begin{tabular}{@{}p{2.0cm}p{1.8cm}p{2.4cm}p{3.0cm}@{}}
  \toprule
  \textbf{Study} & \textbf{Unit} & \textbf{Dependent variable} &
  \textbf{What is missing} \\
  \midrule
  이민영 외 (2026) & Business × week & Page View, Total Clicks &
    Review content and visit context not observed \\[4pt]
  Iselin (1989) & Conceptual & — &
    Defines diversity; not applied to visit context \\[4pt]
  \textcolor{UBIred}{\textbf{This study}} & \textbf{Review × place} &
    \textbf{Revisit, Keyword Count, Rating, Reactions} &
    \textbf{Antecedents of diversity} \\
  \bottomrule
\end{tabular}}
\end{center}
\aside{\scriptsize
  {\color{UBIred}\textbf{The same instrument, one level down.}}
  Naver caps keyword selection at five, so our Keyword Count sits on exactly the
  0--5 scale Lee et al.\ (2026) report.}
\end{frame}

% ════════════════════════════════════════════════════════════════════
\begin{frame}
\frametitle{4. Data \ko{데이터}}
{\small Naver Place reviews, scraped place by place. The unit is the review $i$
at place $j$ — not the business--week.}
\vspace{8pt}
\begin{center}
\begin{tabular}{@{}cccc@{}}
  {\color{UBIred}\textbf{12,537}} & {\color{UBIred}\textbf{7,230}} &
  {\color{UBIred}\textbf{9,843}}  & {\color{UBIred}\textbf{4,148}} \\
  {\scriptsize Reviews} & {\scriptsize Reviewers} & {\scriptsize With text} &
  {\scriptsize With a rating} \\
\end{tabular}
\end{center}
\vspace{4pt}
\aside{\scriptsize
  Missingness is \textbf{structural}, not noise: Naver does not require a star
  rating (66.9\% missing), and a review may be keyword tags with no text at all
  (21.5\%). That is why M3 runs on 4,148 cases.}
\end{frame}

\begin{frame}
\frametitle{4. Variables \ko{변수 소개}}
\begin{center}
{\tiny\color{black}
\begin{tabular}{@{}p{1.4cm}p{2.8cm}p{5.4cm}@{}}
  \toprule
  \textbf{Role} & \textbf{Variable} & \textbf{Definition} \\
  \midrule
  \textcolor{UBIred}{\textbf{Dependent}} \kosub{종속} &
    Revisit \kosub{재방문}\newline Keyword Count \kosub{키워드 수}\newline
    Rating \kosub{별점}\newline Reactions \kosub{반응 수} &
    second or later visit (0/1)\newline
    tags selected, 0--5: information diversity\newline
    star rating, 1--5\newline
    users who reacted: helpfulness \\[2pt]
  \midrule
  \textbf{Independent} \kosub{독립} &
    Companion \kosub{동반자}\newline Purpose \kosub{목적}\newline
    Waited, Reserved\newline Sentiment \kosub{감성} &
    alone, partner, friend, child, family\newline
    travel, outing, daily life, date\newline
    queued; booked in advance (0/1)\newline
    text sentiment, $-1$ to $1$ \\[2pt]
  \midrule
  \textbf{Control} \kosub{통제} &
    Text Length, Images\newline Reviews, Followers\newline
    Review Age, Season\newline Place FE &
    $\ln$ characters; photos attached\newline
    $\ln$ reviewer's counts\newline
    $\ln$ days since posting\newline
    $\delta_j$ absorbs attraction characteristics \\
  \bottomrule
\end{tabular}}
\end{center}
{\tiny\color{black!50}Roles follow 이민영 외 (2026), 소비자학연구 37(1).}
\end{frame}

\begin{frame}
\frametitle{4. Descriptive Statistics \ko{기술통계}}
{\scriptsize Keyword Count is on the same 0--5 scale Lee et al.\ (2026) report.
Note how little the star rating moves — median 5.}
% \tblDescriptives is laid out for the A4 proposal, so it is set at that width
% and then scaled to the slide; typesetting it directly here clips the Max
% column off the right edge.
\begin{center}
\resizebox{\textwidth}{!}{%
  \begin{minipage}{17cm}\color{black}\tblDescriptives\end{minipage}}
\end{center}
\end{frame}

\begin{frame}
\frametitle{4. Models \ko{연구모형}}
{\scriptsize Review-level regressions, place fixed effects $\delta_j$, standard
errors clustered by reviewer.}
{\scriptsize
\begin{align*}
  \text{M1:}\ & \operatorname{logit} P(\textit{Revisit}) = \beta_0 + \mathbf{x}^\top\boldsymbol\beta + \delta_j \\[-3pt]
  \text{M2:}\ & P(\textit{KeywordCount} \le k) = \big[1 + e^{-(\tau_k - \mathbf{x}^\top\boldsymbol\beta - \delta_j)}\big]^{-1} \\[-3pt]
  \text{M3:}\ & P(\textit{Rating} \le k) = \big[1 + e^{-(\tau_k - \mathbf{x}^\top\boldsymbol\beta - \delta_j)}\big]^{-1} \\[-3pt]
  \text{M4:}\ & \textit{Reactions} \sim \text{NegBin}(\mu),\ \
                \ln \mu = \beta_0 + \mathbf{x}^\top\boldsymbol\beta + \gamma\ln(\textit{Age}) + \delta_j
\end{align*}}
\vspace{-6pt}
\aside{\scriptsize
  The link follows the \textbf{measurement scale}: binary (M1), ordinal (M2,
  M3), overdispersed count (M4) — variance $\approx 37\times$ the mean, so
  Poisson is ruled out ($\alpha = 2.53$). All VIF below 2.6.}
\end{frame}

% ════════════════════════════════════════════════════════════════════
\begin{frame}
\frametitle{5. Results \ko{연구결과}}
\begin{center}
{\tiny\color{black}
\begin{tabular}{@{}p{2.0cm}p{2.2cm}p{5.2cm}@{}}
  \toprule
  \textbf{Model} & \textbf{Out-of-sample} &
  \textbf{Leading effects (odds / incidence ratios)} \\
  \midrule
  M1 Revisit \kosub{재방문} & AUC 0.776 &
    reviewer's reviews 1.53\textsuperscript{***}, alone 1.67\textsuperscript{**},
    waited 0.50\textsuperscript{***}, sentiment 0.68\textsuperscript{***} \\[3pt]
  M2 Keyword Count \kosub{키워드 수} & RMSE 1.34 (base 1.86) &
    wrote text 4.39\textsuperscript{***}, visit info 2.18\textsuperscript{***},
    revisitor 1.62\textsuperscript{***} \\[3pt]
  \textcolor{gray}{M3 Rating \kosub{별점}} & \textcolor{gray}{weak, exploratory} &
    \textcolor{gray}{keyword count 1.59\textsuperscript{***}} \\[3pt]
  M4 Reactions \kosub{반응 수} & AUC 0.819 &
    followers 1.96\textsuperscript{***}, text length 1.64\textsuperscript{***},
    \textbf{sentiment 0.76\textsuperscript{**}} \\
  \bottomrule
\end{tabular}}
\end{center}
{\tiny * $p<0.05$, ** $p<0.01$, *** $p<0.001$. Cross-validation grouped by
 reviewer.}
\end{frame}

\begin{frame}
\frametitle{5. Replicates, and New \ko{재현과 발견}}
\aside{\scriptsize
  {\color{UBIred}\textbf{Replicates — sentiment hurts engagement.}}
  More positive reviews get fewer reactions (0.76\textsuperscript{**}).
  Lee et al.\ (2026) find the same sign on clicks ($-0.011$ to $-0.018$).}
\aside{\scriptsize
  {\color{UBIred}\textbf{New — where diversity comes from.}}
  Writing text at all is the strongest predictor of more keywords, then visit
  information, then being a returning visitor. Diversity tracks
  \textbf{engagement}, not the kind of visit.}
{\scriptsize
\begin{itemize}\itemsep2pt
  \item Queueing is the strongest negative signal for returning — and the one a
        manager can act on.
  \item Keyword use has grown over the window ($r = -0.325$ against review age).
\end{itemize}}
\end{frame}

% ════════════════════════════════════════════════════════════════════
\begin{frame}
\frametitle{6. Conclusion \ko{결론}}
{\scriptsize
{\color{UBIred}\textbf{What we find}} \kosub{주요 결과}
\begin{itemize}\itemsep0pt
  \item Diversity tracks how engaged a reviewer is, not the kind of visit.
  \item Sentiment enters engagement negatively, replicating Lee et al.\ (2026).
\end{itemize}
\vspace{6pt}
{\color{UBIred}\textbf{Limitations}} \kosub{한계}
\begin{itemize}\itemsep0pt
  \item Associations only — no keyword-expansion event, no $t-1$ lags.
  \item M1 profiles returning visitors; it does not explain why anyone returns.
\end{itemize}
\vspace{6pt}
{\color{UBIred}\textbf{Next steps}} \kosub{향후 계획}
\begin{itemize}\itemsep0pt
  \item Month fixed effects, to match the two-way specification; Hausman test.
  \item More attractions; validate aspect extraction against manual coding.
\end{itemize}}
\end{frame}

% ── closing ─────────────────────────────────────────────────────────
\begin{frame}
\begin{center}
  \vfill
  {\Large\color{UBIred}\textbf{Thank you}}\par
  \vspace{10pt}
  {\large 감사합니다}\par
  \vspace{18pt}
  {\scriptsize Jorge Adan Soriano Hernandez \quad·\quad
   Department of Data Science, Hoseo University \quad·\quad October 2026}
  \vfill
\end{center}
\end{frame}

\end{document}
